% META: {"id":"1SPE-GEOREP-CO-007","chapitre":"1SPE-GEOMETRIE-REPEREE","type_objet":"corrige","exercice_ref":"1SPE-GEOREP-EX-007","capacites_codes":["C1"],"sources_inspiration":[],"status":"approved","origin":"generated","status_history":["generated","audited","accepted","approved"],"release_acceptance":"RELEASE_OWNER_BATCH_ACCEPTANCE","acceptance_closure_digest":"sha256:9b3ccf9a81c5520fbb7e03b7b2d3e7bf2057a02834b6d908bf3be5f49f3b553f","acceptance_receipt_digest":"sha256:4fad86f0282e0fa4e0419cca1ad6379ae7af5a280c04a4a90880f90a1c044242"}
% BEGIN-VERIFY
% from sympy import *
% x, y = symbols('x y')
% # Droite par A(1,-2) parallele a D: 3x + y - 5 = 0
% # Meme vecteur directeur => meme a,b => 3x + y + c = 0
% # 3*1 + (-2) + c = 0 => c = -1
% assert 3*1 + (-2) - 1 == 0
% END-VERIFY
\begin{corrige}{1SPE-GEOREP-EX-007}

\commentaireMarge{Deux droites parallèles ont le même vecteur directeur, donc les mêmes coefficients $a$ et $b$.}

La droite cherchée a pour équation $3x + y + c = 0$.

$A(1 ; -2)$ est sur cette droite : $3 \times 1 + (-2) + c = 0$, soit $c = -1$.

L'équation est $\boxed{3x + y - 1 = 0}$.

\end{corrige}
